\chapter{Stress}

\section*{Definition}
A state of physiological or psychological strain resulting from demands (stressors) that exceed an individual's perceived coping capacity, eliciting autonomic, endocrine, immune, and behavioral responses.

\section*{Pathophysiology}
Stress activates the sympathetic--adrenal--medullary axis (immediate catecholamine surge) and the hypothalamic--pituitary--adrenal axis (corticotropin-releasing hormone stimulates ACTH, releasing cortisol from the adrenal cortex). Cortisol mobilizes glucose and suppresses non-essential functions; chronic elevation impairs hippocampal neurogenesis, reduces immune competence, and promotes visceral adiposity, hypertension, and insulin resistance.

\section*{Causes}
\begin{itemize}
  \item \textbf{Acute stressors:} Traumatic events (accidents, violence, assault), acute illness, surgery, public speaking, examinations
  \item \textbf{Chronic stressors:} Work-related (burnout, job insecurity, high demand), financial strain, caregiving for ill relative, chronic illness, relationship conflict, discrimination
  \item \textbf{Developmental:} Adverse childhood experiences (abuse, neglect, household dysfunction), academic pressure, social media exposure
  \item \textbf{Environmental:} Noise pollution, overcrowding, unsafe housing, food insecurity, natural disasters
\end{itemize}

\section*{When to See a Doctor}
See your doctor if:
\begin{itemize}
  \item Stress causes functional impairment at work, school, or in relationships
  \item Accompanying suicidal ideation, panic attacks, dissociation, or substance use
  \item Persistent physical symptoms (chest pain, headache, GI distress) despite negative medical workup
\end{itemize}

\section*{Related Symptoms}
\begin{itemize}
  \item Anxiety, irritability, depressed mood, anhedonia, emotional lability
  \item Fatigue, insomnia or hypersomnia, poor concentration, muscle tension, headaches
  \item Palpitations, hyperventilation, diaphoresis, tremor, gastrointestinal upset (nausea, diarrhea)
\end{itemize}
\textbf{Important:} Chronic stress is a transdiagnostic risk factor that amplifies the incidence and severity of cardiovascular, metabolic, psychiatric, and autoimmune disorders.

\section*{Self-Care / Home Management}
\begin{itemize}
  \item Regular aerobic and resistance exercise (150+ minutes per week) reduces cortisol and increases endorphins
  \item Consistent sleep schedule with 7--9 hours per night; limit blue light exposure 1 hour before bed
  \item Mindfulness-based stress reduction (MBSR) or daily meditation practice (10--20 minutes)
  \item Limit caffeine and alcohol; prioritize balanced meals with adequate protein and complex carbohydrates
\end{itemize}

\section*{How Your Doctor Will Evaluate You}
\begin{itemize}
  \item Use validated screening tools: Perceived Stress Scale (PSS-10), PHQ-9 for depression, GAD-7 for anxiety
  \item History to identify specific stressors, coping style, social support, and functional impact
  \item Physical exam and basic labs (CBC, CMP, TSH) to rule out organic mimics of stress symptoms
  \item Assess for secondary conditions: hypertension, insomnia disorders, irritable bowel syndrome, tension-type headache
\end{itemize}

\section*{Treatment Options}
\begin{itemize}
  \item Psychotherapy: CBT, MBSR, acceptance and commitment therapy, or stress inoculation training
  \item Pharmacotherapy for comorbid anxiety or depression: SSRIs/SNRIs (sertraline, escitalopram, venlafaxine)
  \item Biofeedback and relaxation training for somatic symptoms
  \item Social support strengthening: family therapy, support groups, care coordination for high-stress situations
\end{itemize}

\clearpage
